\documentclass[11pt]{article}
\usepackage[margin=1in]{geometry}
\usepackage{amsmath,amssymb,amsthm,booktabs,array,hyperref}
\newtheorem{theorem}{Theorem}
\newtheorem{proposition}[theorem]{Proposition}
\newtheorem{corollary}[theorem]{Corollary}
\newtheorem{remark}[theorem]{Remark}
\title{Spread Obstructions, Frame Colourings, and the Signed Edge Module of \(W(3,3)\)}
\author{Computational manuscript draft}
\date{1 August 2026}

\begin{document}
\maketitle

\begin{abstract}
We study the graph whose 540 vertices are unordered pairs of disjoint totally
isotropic lines of the symplectic generalized quadrangle \(W(3,3)\), with two
vertices adjacent when their canonical cross-matchings share an edge.  The graph
is the edge-disjoint union of 240 cliques of order nine.  A spread determines a
45-vertex maximal independent set which cannot be completed to a resolution;
the obstruction is generated by a fixed-point-free outer involution associated
with a nonsquare symplectic similitude.  For \(q=3,5,7\) the completing-frame
census obeys an exact \(1/q\) law whose support is a perfect matching inside each
spread line.  We separate these finite-geometric results from the representation
theory of the orientation-signed 240-edge module.  That module has one
non-rational constituent, a coexact 90-dimensional real module with Eisenstein
character field and finite integral phase units \(C_6\).  We record explicitly
that proposed electric-charge and homological-flux interpretations of this
\(C_6\) have been withdrawn.  The chromatic decision \(\chi(H)=9\) remains open.
All theorem statements are tagged below as proved, finite-case verified,
likely known with no reference located, computationally open, or withdrawn.
\end{abstract}

\section{Objects and computational conventions}
Let \(W(q,q)\) denote the symplectic generalized quadrangle.  At \(q=3\) it has
40 points, 40 totally isotropic lines, and 240 collinear point pairs.  A
\emph{frame} is an unordered pair of disjoint isotropic lines together with its
canonical four-edge cross-matching.  There are 540 frames.  The frame graph
\(H\) has the frames as vertices, adjacent when their four-edge matchings
intersect.

All finite statements at \(q=3\) were checked from literal point, line, edge,
frame, and group actions.  Statements at \(q=5,7\) are labelled as finite-case
verification unless a uniform argument is supplied.  Code and frozen
certificates are part of the reproducibility record, not substitutes for the
scope labels.

\section{The frame graph}
\begin{proposition}[proved at \(q=3\)]
Each edge of \(W(3,3)\) lies in exactly nine frames.  Those nine frames form a
clique in \(H\), the 240 cliques are edge-disjoint, and
\[
 |V(H)|=540,\qquad |E(H)|=8640,\qquad \deg H=32.
\]
Consequently \(H\) is exactly the union of the 240 edge-indexed \(K_9\)'s.
\end{proposition}

\begin{proof}
A frame containing a fixed edge is obtained by choosing the two disjoint
isotropic lines carrying the endpoints and the compatible cross-matching.
Literal enumeration gives nine choices.  Each frame has four matching edges,
and for each such edge there are eight other frames in its clique, giving
degree \(4\cdot8=32\).  The clique edge count is
\(240\binom92=8640\), equal to \(540\cdot32/2\).  The equality, together with
the direct check that adjacent frames share one matching edge, proves that no
graph edge is counted twice.
\end{proof}

\begin{corollary}
The assertion \(\chi(H)=9\) is equivalent to assigning nine colours to the 540
frames so that every edge-indexed \(K_9\) contains every colour exactly once.
\end{corollary}

\section{Spread traps}
A spread \(S\) consists of \(q^2+1\) pairwise disjoint isotropic lines
partitioning the points.

\begin{theorem}[proved uniformly at the incidence level]
The \(\binom{q^2+1}{2}\) pairs of lines in a spread are frames and form an
independent set of \(H\).  Their matchings cover precisely the collinear point
pairs not lying on the spread lines.  The uncovered set has size
\[
 (q^2+1)\binom{q+1}{2}.
\]
\end{theorem}

\begin{theorem}[verified for all 36 spreads at \(q=3\)]
The 45 frames determined by a spread form a maximal independent set, although
\(\alpha(H)=60\).  The 15 candidate frames supported on the 60 uncovered edges
touch only 20 of those edges, leaving 40 edges in no candidate frame.
Therefore no spread independent set extends to a 60-frame resolution class.
\end{theorem}

\section{The \(1/q\) support law}
\begin{theorem}[finite-case verified for \(q=3,5,7\)]
The number of admissible completing frames is
\[
 \frac{q(q^2+1)}2.
\]
Their matching edges occupy
\[
 \frac{(q+1)(q^2+1)}2
\]
distinct uncovered edges, exactly a fraction \(1/q\) of the uncovered set, and
every occupied edge has multiplicity \(q\).
\end{theorem}

The measured ratios are \(20/60\), \(78/390\), and \(200/1400\).
The support consists of a perfect matching of size \((q+1)/2\) inside each
spread line.  This explains the odd-\(q\) branch.  At \(q=2\), where a spread
line has odd cardinality, the measured candidate set is empty.  A uniform
written proof that every candidate support must be this perfect matching remains
an open exposition task.

\section{The spread involution}
\begin{theorem}[existence uniform for odd \(q\); uniqueness verified at \(q=3\)]
Let \(\mu\) be a nonsquare and let \(g\in GSp(4,q)\) satisfy \(g^2=\mu I\).
Then the induced projective transformation is fixed-point-free and preserves
the associated spread linewise.  At \(q=3\), each of the 36 spreads has a unique
nontrivial linewise stabilizer \(\sigma_S\), it is central in the spread
stabilizer of order 1440, and the candidate matching edges are its 2-cycles.
\end{theorem}

A fixed projective point would require an eigenvalue \(\lambda\in\mathbb F_q\)
with \(\lambda^2=\mu\), impossible for nonsquare \(\mu\).  At \(q=3\), the 72
similitudes of this type give 36 projective involutions.  The square-multiplier
branch gives 270 inner involutions.  The framework is standard; no reference
for this exact \(36/270\) split has been located, so it should be treated as
likely known rather than advertised as new.

\section{The orientation-signed edge module}
Over \(PGSp(4,3)\), the signed 240-edge module decomposes multiplicity-freely as
\[
 V=15\oplus24\mid81\mid30\oplus90,
\]
with exact, harmonic, and coexact dimensions \(39,81,120\).
The exact \(15\oplus24\) is inherited from the real point-permutation module.
The coexact 90 restricts to a complex-conjugate \(45\oplus\overline{45}\) pair
for \(PSp(4,3)\), and
\[
 \operatorname{End}_{PSp}(V)\cong\mathbb R^4\times\mathbb C.
\]
The 90 is the only non-rational block.  Its character field is
\(\mathbb Q(\omega)\), and the finite integral phase units are \(C_6\).

The exact structural interpretation is limited: the characteristic \(C_3\)
inside the finite centralizer torsion acts trivially on the rational
150-dimensional block sum and faithfully on the Eisenstein 90.  The outer
involution acts by inversion.  Thus the phase is an internal cyclotomic sector
marker reversed by chirality.

\begin{table}[h]
\centering
\begin{tabular}{p{0.34\linewidth}p{0.56\linewidth}}
\toprule
Proposed reading & Status and falsifier\\
\midrule
Electric charge from a Gauss-law sector &
Withdrawn: the phase lies in the coexact 90, not the exact source/divergence
block.\\
Homological or Dirac flux quantum &
Withdrawn: the clique complex is torsion-free at the tested integral boundary
maps; no \(\mathbb Z/3\) or \(\mathbb Z/6\) torsion supports such a quantum.\\
QCD colour, generation number, neutrino assignment &
Not derived; no such identification is claimed.\\
Internal cyclotomic sector marker &
Supported representation-theoretically: unique odd torsion of the integral
centralizer, supported exactly on the 90 and inverted by the outer involution.\\
\bottomrule
\end{tabular}
\end{table}

\section{The unresolved colouring problem}
The strongest certified formulations currently give only bounded
\texttt{UNKNOWN} results.  Spread-count branching reduced the historical
search to 60,909 CP-SAT branches.  Audited geometric symmetry cuts reduce a
different formulation but do not close it.  A combined spread-signature
formulation with all 40 point-transvection orbit cuts is exact and removes
96.8866\% of the 25,920 geometric images of a known proper 14-colouring, yet a
bounded nine-colour MILP still returns \texttt{UNKNOWN}.  No chromatic theorem
is inferred from these runs.

\section{Ownership, prior art, and novelty boundary}
The general facts about non-real symplectic characters and twisted
Frobenius--Schur indicators belong to the literature represented in the
repository by Gow and Vinroot.  The Weil-pair chirality frame was already in
the repository's Passes 353 and 355; earlier rederivations are not claimed as
novel.  The \(E_8\) Coxeter \(C^5\) six-cycle fibration was already established
in BT1745--BT1751 and is inequivalent to the full 240-dimensional internal
endomorphism \(C_6\) module.

The paper-level contribution claimed here is narrower: the exact frame-graph
decomposition; the spread trap and finite \(1/q\) support census; the literal
\(\sigma_S\) obstruction at \(q=3\); and the signed-edge-module localization
with explicit retraction boundaries.

\section{Reproducibility checklist}
\begin{center}
\begin{tabular}{lll}
\toprule
Claim & Evidence type & Status\\
\midrule
240 edge-disjoint \(K_9\)'s & literal incidence census & exact\\
36 spread traps at \(q=3\) & exhaustive finite enumeration & exact\\
\(1/q\) law & \(q=3,5,7\) enumerations & finite-case verified\\
\(\sigma_S\) uniqueness & full \(q=3\) group enumeration & exact at \(q=3\)\\
nonsquare-similitude existence & algebraic argument & odd \(q\)\\
signed-module decomposition & exact projectors/characters & exact\\
internal \(C_6\) & character field and integral units & exact\\
\(\chi(H)=9\) & bounded solver runs & open\\
\bottomrule
\end{tabular}
\end{center}

\section{Open problems}
\begin{enumerate}
\item Decide whether \(\chi(H)=9\).
\item Prove or refute that the maximum spread intersection of a resolution class
is 13; 13 is attained and 14 remains undecided.
\item Give a uniform proof of the within-line perfect-matching mechanism.
\item Locate prior art for the \(36/270\) nonsquare/square multiplier split.
\item Determine whether the internal cyclotomic sector marker has any canonical
role beyond representation theory, without importing a physical label.
\end{enumerate}

\end{document}
